V rámci tohto tímového projektu sme úspešne navrhli, implementovali a 
otestovali komplexný vnorený riadiaci systém pre model robotického 
manipulátora so šiestimi stupňami voľnosti (6-DOF). Podarilo sa nám 
naplniť všetky hlavné body definované v zadaní a priniesť plne funkčné 
mechatronické riešenie postavené na hardvérovej platforme mikrokontroléra 
STM32. Celé programové vybavenie bolo implementované v jazyku C++ s využitím
knižnice Eigen bez 
použitia rozsiahlych frameworkov pre robotiku (ako napr. ROS), čo nám 
poskytlo oveľa hlbší pohľad do problematiky tohto riešenia.

Architektonickým základom bolo využitie operačného systému reálneho času 
(FreeRTOS). Determinizmus systému sme dosiahli vďaka dôslednému rozdeleniu
 procesov na samostatné paralelné úlohy s vhodnou prioritizáciou
 , ktoré bezpečne a s nízkou 
 latenciou zvládajú výpočet kinematiky, generovanie riadiacich signálov 
 pre servomotory aj komunikáciu naprieč procesmi prostredníctvom správ a 
 frontov.

Jadrom práce bolo zvládnutie problematiky 
priestorovej kinematiky. Návrh numerického riešenia inverznej kinematickej 
úlohy (IKÚ) na báze Newtonovej metódy a maticových operácií z knižnice 
Eigen bol úspešne aplikovaný do výpočtov v reálnom čase.

Ďalším dôležitým splneným míľnikom z pohľadu zadania a hardvérovej 
integrácie bolo spojazdnenie prehľadného používateľského rozhrania, 
ktoré kombinuje LCD displej a joystick. Jeho prostredníctvom bola 
zabezpečená spätná väzba a zadávanie cieľových priestorových 
súradníc koncového efektora. Navyše bol systém rožšírený o aplikáciu 
3D mapovania priestoru integráciou laserového ToF senzora (VL53L1X), na 
ktorú nadväzuje funkčné logovanie vyhodnotených dát priamo na SD kartu.
Taktiež sme si vyškúšali skonštruovať ďalší model manipulátora na 
báze 3D tlače a mechanicky navrhnúť jeho vybranú 
časť manipulátora, ktorá viedla k vytvoreniu nového manipulátora.

Zhodnotením celej našej práce konštatujeme, že tento tímový projekt splnil 
všetky definované ciele integrácie pokročilého algoritmu a embedded 
hardvéru. Aplikovaním teoretických matematických modelov z oblasti robotiky 
priamo do softvéru s limitovaným hardvérovým výkonom sme získali cenné 
praktické skúsenosti. Náš systém je robustný, preukázal funkčnosť aj 
napriek mechanickým obmedzeniam 6-DOF modelu a zároveň tvorí otvorený 
systém, ktorý je plne pripravený pre možné budúce rozsiahlejšie experimenty 
a vylepšenia.